\documentclass[UTF8,a4paper,11pt]{ctexart}
\usepackage[margin=2cm]{geometry}
\usepackage{amsmath,amssymb,fancyhdr,hyperref}
\hypersetup{hidelinks}
\pagestyle{fancy}\fancyhf{}\fancyfoot[C]{\thepage}\renewcommand{\headrulewidth}{0pt}
\setlength{\parindent}{0pt}\setlength{\parskip}{5pt}
\begin{document}
\begin{center}{\LARGE 有来源版：教师解析与评分}\end{center}
5题共100分，建议60分钟。第1、2题检查概念，第3、4题训练基本不等式，第5题为新定义挑战。来源、页码和改编范围见第3页。以下解析重新推导，分值为本卷分值，不是原卷分值；其他正确解法按对应推理给分。

\textbf{1．答案：$a=1$ 或 $a=2$。}

恰有2个子集意味着 $A$ 恰有一个元素。理由是空集只有1个子集；单元素集合有2个子集；若有两个不同元素，至少能列出4个不同子集。

当 $a=1$ 时，方程成为 $-2x+1=0$，故 $A=\{\frac12\}$，符合要求。
当 $a\ne1$ 时，方程为一元二次方程，恰有一个不同实根要求
\[\Delta=(-2)^2-4(a-1)=8-4a=0,\]
故 $a=2$，此时 $(x-1)^2=0$，$A=\{1\}$，符合要求。
因此 $a=1$ 时 $A=\{\frac12\}$；$a=2$ 时 $A=\{1\}$。

\textbf{评分（16分）：}子集数转为一个元素4分；$a=1$ 的情形及集合4分；$a\ne1$ 的判别式与参数4分；相应集合与完整结论4分。

\textbf{2．答案：充分不必要条件。}

若 $a+b\le4$，由正数基本不等式可得
\[2\sqrt{ab}\le a+b\le4,\quad\text{故}\quad ab\le4.\]
这证明了充分性。反之，取 $a=4,b=1$，有 $ab=4\le4$，但 $a+b=5>4$，所以 $q$ 不能推出 $p$，必要性不成立。

\textbf{评分（16分）：}充分性证明6分；满足 $q$ 且不满足 $p$ 的有效反例6分；正确判断4分。单独举例不能替代充分性的普遍证明。

\textbf{3．答案：最小值 $2\sqrt2$，在 $a=b=\sqrt2$ 时取得。}

因为各项为正，连续使用基本不等式，得
\[\frac1a+\frac a{b^2}+b\ge2\sqrt{\frac1a\cdot\frac a{b^2}}+b
=\frac2b+b\ge2\sqrt2.\]
第一步取等要求 $\frac1a=\frac a{b^2}$，由正数条件等价于 $a=b$；第二步取等要求 $\frac2b=b$，即 $b=\sqrt2$。两处可同时取等，且只有 $a=b=\sqrt2$ 时同时取等。

独立复核：
\[\frac1a+\frac a{b^2}+b-2\sqrt2
=\left(\frac1{\sqrt a}-\frac{\sqrt a}{b}\right)^2
+\left(\sqrt b-\sqrt{\frac2b}\right)^2\ge0.\]

\textbf{评分（20分）：}第一步估计6分；第二步估计6分；两个等号条件各3分；最小值及可达性结论2分。
\newpage
\textbf{4．证明。}

（1）由 $(a+b+c)^2=0$ 得
\[ab+bc+ca=-\frac12(a^2+b^2+c^2).\]
因为 $abc=1$，三个数均不为0，所以平方和严格大于0，所求式严格小于0。

（2）先判断符号。由乘积为正且三个数非零，负数的个数只能为0或2；若没有负数，和必为正，与题设矛盾。因此恰有一个正数和两个负数。

不妨记最大者为 $a$，则 $a>0,b<0,c<0$，且 $a=-(b+c)$、$bc=1/a$。由 $(b-c)^2\ge0$，得 $(b+c)^2\ge4bc$，因此
\[a^3=a(b+c)^2\ge4abc=4.\]
因为 $a>0$，故 $a\ge\sqrt[3]{4}$，即 $M\ge\sqrt[3]{4}$。也可对正数 $-b,-c$ 使用基本不等式。

取等核查：令 $t=\sqrt[3]{\frac12}$，取 $(a,b,c)=(2t,-t,-t)$，则和为0、积为1，最大者 $2t=\sqrt[3]{4}$，证明界可达到。

\textbf{评分（24分）：}第（1）问8分，展开关系4分、解释严格性4分；第（2）问16分，符号分析6分、选最大者并转化4分、平方非负或基本不等式4分、结论2分。原设问不要求求等号条件，因此未写取等核查不扣分。

\textbf{5．答案与证明。}

（1）前两项恒为1，第三项为 $y_3$，所以 $\alpha*\beta=2+y_3$。要求等于3，等价于 $y_3=1$，而 $y_1,y_2$ 各可取0或1。因此全部解为
\[(0,0,1),\quad(0,1,1),\quad(1,0,1),\quad(1,1,1).\]

（2）逐项观察：$x_i+y_i-x_iy_i$ 在 $(x_i,y_i)=(0,0)$ 时为0，在其余三种情形下均为1。因此 $\alpha*\beta$ 等于“至少有一个1”的位置数。

设 $\alpha$ 有 $r$ 个1，$\beta$ 有 $s$ 个1，两者同为1的位置有 $t$ 个。因为 $x_i^2=x_i$，有 $\alpha*\alpha=r$，同理 $\beta*\beta=s$。条件即 $r+s=n$，所以
\[\alpha*\beta=r+s-t=n-t,\qquad 0\le t\le\min(r,s).\]
最大值为 $n$：让 $\alpha$ 全为1、$\beta$ 全为0，即可取得。

若 $n=2m$，必有 $\min(r,s)\le m$，故 $n-t\ge m$。取两个数组相同，均在前 $m$ 个位置为1、其余为0，则 $r=s=t=m$，取得最小值 $m=n/2$。

若 $n=2m+1$，因为 $r,s$ 是整数且和为 $2m+1$，至少一个不超过 $m$，故 $t\le m$，$n-t\ge m+1$。取 $\alpha$ 的前 $m$ 位为1，$\beta$ 的前 $m+1$ 位为1，其余全为0，便有 $r=m,s=m+1,t=m$，取得最小值 $m+1=(n+1)/2$。

\textbf{评分（24分）：}第（1）问8分，得到 $y_3=1$ 得4分、完整列举4分；第（2）问16分，位置计数解释及 $n-t$ 得4分、最大值及构造4分、最小值的界4分、奇偶两种取等构造4分。题设没有要求 $\alpha\ne\beta$，偶数情形可以取相同数组。
\newpage
\textbf{来源与改编记录}

以下均已对照公开PDF的题目页面核查；公开件由资料网站转载或整理，不将其称为学校／考试机构官方发布版本。完整链接保存在同目录的\texttt{来源与核查.md}中。链接编号与下表相同。

\textbf{[S1]} 2024年10月北京景山学校高一上学期期中数学试卷，公开PDF共13页，由北京中考在线网站提供下载。

第1题取自原第12题（PDF第2页）：方程与子集数条件不变，填空改为解答，新增写出对应集合的要求。

第2题取自原第7题（PDF第1页）：正数条件和两个不等式不变，去掉选项，改为判断并论证。

第5题取自原第21题（I）（II）（PDF第3页）：保留运算定义及这两问的数学条件，以数组解释原有序数组记号，新增一个记号示例并明确要求验证最值可达；不选用涉及更强组合推理的第（III）问。

\textbf{[S2]} 2021年天津高考数学第13题，公开《2021年天津市高考数学试卷》PDF第2页（共20页）。

第3题保留原表达式与正数条件，填空改为证明最小值并写出取等点。\href{https://17094505.s21i.faiusr.com/61/ABUIABA9GAAg2bzahwYovJaMRQ.pdf}{点击查看公开PDF}。

\textbf{[S3]} 2020年全国III卷文科数学第23题，公开《2020年高考全国III卷真题及答案》PDF第6页（共10页），自主选拔在线转载。

第4题保留全部条件与两项待证结论，仅用 $M$ 代替原题 $\max\{a,b,c\}$ 并调整措辞。\href{https://cdn.zizzs.com/zixunzhan/202310/d89f23fd-c957-4f40-b6ef-86aa0eb869f0.pdf}{点击查看公开PDF}。

\textbf{教学目标、易错点与评价依据}

第1题（中等）：子集数、实根个数和参数分类。预期先判断集合元素数，再讨论是否为二次方程；重点检查是否遗漏 $a=1$，以及是否把重根当两个集合元素。

第2题（基础向中等过渡）：充分必要条件与基本不等式的联用。预期普遍证明一个方向、用反例否定反向；重点检查是否把举例当作充分性证明。

第3题（中等偏难）：连续使用两次基本不等式。预期先处理前两项，再处理与 $b$ 的和；重点检查两处取等条件能否同时实现。

第4题（较难）：在混合正负数中识别可用不等式的对象。预期先判断符号，再用平方非负或对 $-b,-c$ 使用基本不等式；重点检查严格不等号的依据，以及乘不等式前是否知道乘数为正。

第5题（挑战）：读懂新定义，将代数运算解释为集合中位置的计数。预期四种取值列表、重复位置计数、奇偶分类并构造；重点检查下界与可达性是否均说明。无需数列公式、排列组合公式或向量知识。

难度为针对开学三周学生的设计判断，未经学生实测。第1题使用初中一元二次方程判别式，第4题使用立方根，第5题使用奇偶分类；如学生卡在记号，可在讲评时用三格或四格的0、1数组示范。学生卷为2页A4，100\%实际大小、长边翻转双面打印。
\end{document}
